\documentclass{article}

\usepackage[utf8]{inputenc}
\usepackage[T1]{fontenc}
\usepackage{helvet}
\usepackage{geometry}
\usepackage{xcolor}
\usepackage{eso-pic}
\usepackage{fancyhdr}
\usepackage{parskip}
\usepackage{microtype}
\usepackage[english, ukrainian]{babel}
\usepackage{hyperref}
\usepackage{titlesec}
\usepackage{tocloft}

\definecolor{nato}{HTML}{293757}
\definecolor{footergray}{gray}{0.65}

\geometry{a4paper,left=3cm,right=3cm,top=3cm,bottom=3cm}
\renewcommand{\familydefault}{\sfdefault}
\setcounter{secnumdepth}{3}

\titleformat{\section}{\color{nato}\Huge\bfseries}{\thesection.}{1em}{}
\titlespacing*{\section}{0pt}{30pt}{12pt}

\titleformat{\subsection}{\color{nato}\LARGE\bfseries}{\thesubsection.}{1em}{}
\titlespacing*{\subsection}{0pt}{20pt}{8pt}

\fancyhf{}
\fancyhead[R]{\small\color{footergray} 31 березня 2026}
\fancyfoot[L]{\small\color{footergray} КЗЗІ. Модель профілів безпеки.}
\fancyfoot[R]{\small\color{footergray}\thepage}

\newcommand\HUGE[1]{\fontsize{40}{40}\selectfont #1}

\renewcommand{\headrulewidth}{0pt}
\renewcommand{\footrulewidth}{0.4pt}
\renewcommand{\footrule}{\color{footergray}\hrule width \textwidth height 0.4pt}

\begin{document}

\pagestyle{fancy}

\begin{titlepage}
  \AddToShipoutPictureBG*{\AtPageLowerLeft{\color{nato}\rule{\paperwidth}{\paperheight}}}
  \color{white}\thispagestyle{empty}\vspace*{4cm}\centering
  {\Huge  \textbf{Комплексна система захисту iнформацiї} \par} \vspace{2cm}
  {\HUGE  \textbf{Модель профілів безпеки} \par} \vspace{2.5cm}
  \vfill
  {\large 2026 \copyright\ Інформаційні Судові Системи \par}
\end{titlepage}

\newpage

\tableofcontents

\section{Вступ}
На підприємствах (недержавних суб’єктах) у 2026 році КЗЗІ (Комплексна система захисту інформації) переходить на модель профілів безпеки відповідно до Постанови КМУ № 712 від 18 червня 2025 року.

\section{Нормативно-правова база (НПА)}
\begin{itemize}
    \item \textbf{Постанова КМУ № 712 від 18.06.2025} («Порядок розроблення та затвердження профілів безпеки...» та «Порядок авторизації...»).
    \item \textbf{Закон України «Про судоустрій і статус суддів»} \href{https://zakon.rada.gov.ua/laws/show/1402-19}{(№ 1402-VIII від 02.06.2016)}.
    \item \textbf{Закон України «Про захист інформації в інформаційно-комунікаційних системах»}.
    \item \textbf{Закон України «Про основні засади забезпечення кібербезпеки України»}.
    \item \textbf{Нормативні документи ТЗІ} (НД ТЗІ 3.6-004-21, 3.6-005-21, 3.6-006-24 тощо).
\end{itemize}

\section{Основні типи профілів безпеки}
\begin{itemize}
    \item \textbf{Базовий профіль} — затверджується Адміністрацією Держспецзв’язку (наказом).
    \item \textbf{Галузевий профіль} — розробляється галузевим органом, погоджується з Держспецзв’язку та затверджується наказом/рішенням відповідного органу.
    \item \textbf{Цільовий профіль безпеки (ЦПБ)} — індивідуальний для конкретної системи підприємства. Саме на його основі створюється/модернізується КЗЗІ.
\end{itemize}

\section{Формат затвердження та прийняття на підприємстві}
Для цільового профілю безпеки на підприємстві (власник/розпорядник системи) затвердження відбувається внутрішнім організаційно-розпорядчим документом:
\begin{itemize}
    \item Наказом керівника підприємства (директора, генерального директора тощо).
    \item Або рішенням уповноваженого органу управління (для ТОВ, АТ тощо — протоколом зборів засновників/наглядової ради, якщо це передбачено статутом).
\end{itemize}

Цільовий профіль розробляється на основі базового (або галузевого, якщо є) з урахуванням:
\begin{itemize}
    \item Результатів оцінки ризиків.
    \item Особливостей системи.
    \item Нормативних документів ТЗІ (НД ТЗІ 3.6-004-21, 3.6-005-21, 3.6-006-24 тощо).
\end{itemize}

\section{Процес (коротко)}
\begin{enumerate}
    \item Обстеження системи та оцінка ризиків.
    \item Формування ЦПБ (сукупність заходів захисту).
    \item Затвердження внутрішнім наказом керівника підприємства.
    \item Документування (політика безпеки інформації тощо).
    \item Впровадження заходів, оцінка відповідності.
\end{enumerate}

Цільовий профіль не потребує зовнішнього затвердження Держспецзв’язку для звичайних підприємств (на відміну від базових/галузевих). Він є основою для декларування/авторизації відповідності системи захисту.

\newpage
\section{Ієрархія профілів безпеки}
Відповідно до \href{https://zakon.rada.gov.ua/laws/show/1402-19}{Закону України «Про судоустрій і статус суддів»}, стаття 17 «Система судоустрою»:
\begin{enumerate}
    \item Судоустрій будується за принципами територіальності, спеціалізації та інстанційності.
    \item Найвищим судом у системі судоустрою є Верховний Суд.
    \item Систему судоустрою складають: 1) місцеві суди; 2) апеляційні суди; 3) Верховний Суд.
\end{enumerate}

На базі цього визначено таксономію профілів безпеки за відповідними OID:
\begin{itemize}
    \item \textbf{1.2.804.3.1.2.1} — Базовий профіль безпеки (L1)
    \item \textbf{1.2.804.3.1.2.2} — Галузевий профіль безпеки (L2)
    \item \textbf{1.2.804.3.1.2.3} — Цільовий профіль безпеки вищих судів (L3)
    \begin{itemize}
        \item 1.2.804.3.1.2.3.1 — Велика Палата Верховного Суду
        \item 1.2.804.3.1.2.3.2 — Касаційний адміністративний суд
        \item 1.2.804.3.1.2.3.3 — Касаційний господарський суд
        \item 1.2.804.3.1.2.3.4 — Касаційний кримінальний суд
        \item 1.2.804.3.1.2.3.5 — Касаційний цивільний суд
    \end{itemize}
    \item \textbf{1.2.804.3.1.2.4} — Цільовий профіль безпеки вищих спеціалізованих судів (L3)
    \begin{itemize}
        \item 1.2.804.3.1.2.4.1 — Вищий суд з питань інтелектуальної власності
        \item 1.2.804.3.1.2.4.2 — Вищий антикорупційний суд
        \item 1.2.804.3.1.2.4.3 — Спеціалізований окружний адміністративний суд
        \item 1.2.804.3.1.2.4.4 — Спеціалізований апеляційний адміністративний суд
    \end{itemize}
    \item \textbf{1.2.804.3.1.2.5} — Цільовий профіль судів (L3)
    \begin{itemize}
        \item 1.2.804.3.1.2.5.1 — Цільовий профіль безпеки місцевих судів (L4)
        \item 1.2.804.3.1.2.5.2 — Цільовий профіль безпек апеляційних судів (L4)
    \end{itemize}
    \item \textbf{1.2.804.3.1.2.6} — Цільовий профіль безпеки органів та установ у системі правосуддя (L3)
    \begin{itemize}
        \item 1.2.804.3.1.2.6.1 — ДСА (L4)
        \item 1.2.804.3.1.2.6.2 — ТУ ДСА, допоміжні установи (L4)
        \item 1.2.804.3.1.2.6.3 — Рада суддів України (L4)
        \item 1.2.804.3.1.2.6.4 — ВРП (L4)
        \item 1.2.804.3.1.2.6.5 — ВККСУ (L4)
        \item 1.2.804.3.1.2.6.6 — Національна школа суддів України (L4)
        \item 1.2.804.3.1.2.6.7 — ГРД та ГРМЕ (L4)
        \item 1.2.804.3.1.2.6.8 — Служба судової охорони (L4)
    \end{itemize}
\end{itemize}


Якщо у вас державна система, об’єкт КІІ чи специфічна галузь — порядок може відрізнятися (погодження/затвердження
з Держспецзв’язку або галузевим органом). Рекомендую звернутися до відділу ТЗІ або спеціалізованої організації для конкретної системи.

\end{document}
